\maketitle
\begin{abstract}
This report documents the design, implementation, and evaluation of eight natural
language processing models built entirely from scratch using PyTorch. The work is
divided into two tasks: (1) text generation using Markov Chains, vanilla RNNs,
and LSTMs at both character and word level, and (2) question-answering chatbots
using an LSTM encoder-decoder with Bahdanau attention and a full Transformer
encoder-decoder with multi-head self-attention. No pretrained weights or external
NLP libraries were used.

Results show that the character-level LSTM achieves the lowest cross-entropy loss
among all text-generation models (0.2943 vs.\ 0.6689 for the RNN), a 56\%
reduction attributable to gating mechanisms that prevent vanishing gradients.
In the chatbot task, the LSTM encoder-decoder converges to a final loss of 0.1014
and produces factually correct answers on 10 test questions, while the Transformer
chatbot (final loss 4.3724) exhibits repetition collapse due to insufficient
training data and the conservative warmup learning-rate schedule. The Transformer
trains 3× faster per epoch owing to full sequence parallelism, but requires
significantly more data and epochs to converge on this small dataset. These results
illustrate that architecture choice, training strategy, and dataset size interact
in non-trivial ways — and that the architecturally superior model is not always
the practically superior one under constrained conditions.
\end{abstract}

\clearpage
